\documentclass[11pt]{article}
\usepackage[margin=1in]{geometry}
\usepackage{graphicx}
\usepackage{booktabs}
\usepackage{amsmath}
\usepackage{hyperref}
\usepackage{natbib}

\title{Self-Falsifying Adaptive Systems: An Executable Falsification Workflow for AI-Assisted Research}

\author{Physis Research Program}
\date{\today}

\begin{document}
\maketitle

\begin{abstract}
We describe an executable falsification workflow implemented in a
production software system: claims are recorded with their evidence,
experiments are run with construction-matched nulls, verdicts are
assigned, revisions are tracked, and failed branches are retained as
research evidence. The workflow was stress-tested over three months of
autonomous AI-assisted research, producing both settled results
(cross-embedder ontology agreement near chance) and a catalog of
twelve failure modes where the null caught the error. We argue that the
falsification workflow, not the grand architecture, is the transferable
contribution.
\end{abstract}

\section{Introduction}

Autonomous research agents (Sakana AI Scientist, Robot Scientist, and
others) typically follow a loop: generate hypotheses, run experiments,
interpret results, update beliefs. The loop is usually described as
hypothesis-driven science. What is less often described is the
infrastructure that makes the loop honest: how are hypotheses recorded?
How are experiments designed to be falsifiable? How are negative results
preserved?

We report from a system that implements this infrastructure explicitly.
Every claim has a ledger entry with a verdict (Supported, Contradicted,
Partially Supported, Underpowered, Artifact, Superseded). Every
experiment specifies a construction-matched null before running. Every
correction is recorded and linked to the original claim. Failed branches
are never deleted.

The system is Physis, a 200+ module Rust application for ontology
management, coherence tracking, and agent coordination. Over three
months of development, the system ran 90+ experiments across five
research tracks. The experiments produced settled results --- including a
demonstration that cross-embedder ontology agreement is near chance
(ARI ~0.10) and that cross-agent failure priors transfer with a
coverage gate --- and a dozen failure modes where the null caught the
error.

\section{Architecture}

The falsification workflow has five stages:

\begin{enumerate}
\item \textbf{Claim registration.} A hypothesis is recorded with its
predicted effect, the experiment that will test it, and the null that
could kill it. The claim has a unique ID and a timestamp.

\item \textbf{Experiment execution.} The experiment is run. The null is
run in the same execution, on the same data, with the same code path.
The null is not an afterthought; it is a required input.

\item \textbf{Verdict assignment.} The system assigns a verdict:
Supported (effect size excludes the null), Contradicted (the null
explains the result), Partially Supported (the effect survives some
controls but not others), Underpowered (the experiment cannot resolve
the question), Artifact (the result is explained by a confound), or
Superseded (a later experiment modifies the claim).

\item \textbf{Revision tracking.} When a claim is revised, the original
is marked Superseded, not deleted. The revision chain is preserved:
claim -> experiment -> verdict -> revision -> new claim.

\item \textbf{Failed branch retention.} Experiments that produce null
results are retained in the repository. Failed code is marked
deprecated, not deleted. The experimental trail is never erased.
\end{enumerate}

\section{Stress Test: What the Workflow Caught}

The twelve failure modes cataloged in the companion paper (Paper 5) were
all caught by the workflow:

\begin{itemize}
\item \textbf{E14 (wrong injection distribution).} The claim ``geometry
scores misfilings at AUC 0.880'' was registered with a null: injection
at the nearest alternative cell. The null produced AUC 0.506, at chance.
The claim was marked Contradicted. Without the pre-registered null, the
0.880 would have been reported as a positive.

\item \textbf{E83 (generated fixture circularity).} The claim ``the
antonym rule detects contradiction'' was registered with a null:
hand-written fixtures without lexical antonyms. The null produced
recall 0.143. The claim was marked Artifact.

\item \textbf{E85 (adversarial collapse).} The claim ``structured rules
achieve precision 1.000'' was registered with a null: adversarial
fixtures designed to fire each rule wrongly. The null produced
precision 0.000. The claim was marked Contradicted.

\item \textbf{E17 (confound check that stops early).} The claim
``ANCHORVOC is not circular'' was registered with three tests. The
workflow required all three to pass. Test 3 failed, and the claim was
marked Artifact, even though Tests 1 and 2 passed. Without the
requirement to test all three, the script's first verdict (which read
only Test 1) would have been accepted.

\item \textbf{E10 (pre-flight overridden).} The power analysis returned
STOP, and the run proceeded anyway. The workflow records this as a
violation of the pre-registration contract. The claim is marked
Underpowered.
\end{itemize}

In each case, the workflow caught the error because the null was
pre-registered, the baseline was rerun in the same execution, and the
verdict was assigned by the system, not by the agent.

\section{Comparison with Existing Systems}

\textbf{Sakana AI Scientist} (2024, 2025) implements an autonomous
research loop: generate hypotheses, run experiments, interpret results,
update beliefs. The loop is hypothesis-driven, but the falsification
infrastructure (how hypotheses are recorded, how experiments are designed
to be falsifiable, how negative results are preserved) is not described
in detail. Our contribution is the infrastructure, not the loop.

\textbf{AGM belief revision} (Alchourron, Gardenfors \& Makinson, 1985)
provides a formal theory of belief change: contraction, expansion,
revision. Our workflow is an implementation of these principles in a
running system, with the addition of construction-matched nulls and
machine-readable experiment histories.

\textbf{Doyle's TMS} (1979) is a truth maintenance system: beliefs are
justified, and a belief lives while a justification holds. Our workflow
is similar in spirit but operates on empirical claims (with nulls and
effect sizes) rather than logical propositions.

\textbf{Robot Scientist} (King et al., 2009) generates hypotheses,
runs experiments, and interprets results. The Robot Scientist's
hypotheses are formalized in logic; our claims are formalized as
hypothesis + null + verdict + revision chain.

The distinctive feature of our workflow is the \textbf{retained failed
branch}: the experiment ledger preserves not just the surviving claims
but the full chain of claims that were tested and failed. This makes the
experimental trail reproducible and auditable.

\section{Limitations}

\textbf{Single system, single domain.} The workflow was tested on one
system (Physis) in one domain (ontology management and agent
coordination). Its generalizability to other domains is untested.

\textbf{No comparison group.} We cannot claim that the workflow
produces better science than a workflow without it. Such a comparison
would require a controlled study with two research groups using
different workflows.

\textbf{Agent compliance.} The workflow catches errors only if the agent
follows it. E10 (pre-flight overridden) shows that an agent can bypass
the workflow. The workflow records the violation but cannot prevent it.

\textbf{Scale.} The workflow was tested on 90+ experiments over three
months. Its performance at larger scales (hundreds of agents, thousands
of experiments) is unknown.

\section{Conclusion}

The falsification workflow --- claim registration, pre-registered nulls,
in-run baselines, revision tracking, failed branch retention --- is a
transferable contribution. It is not tied to the Physis architecture or
to any particular research domain. Any system that uses AI agents to
design, implement, and interpret experiments can adopt it.

The twelve failure modes cataloged in the companion paper demonstrate
that the workflow is not ceremonial: in each case, the null caught an
error that would otherwise have been reported as a positive result.

\end{document}
